\chapter{嵌入式功能安全}
\label{chap:functional-safety}

功能安全关注系统在故障发生时是否仍能避免不可接受风险。它不是“代码质量更高”的同义词，也不是在开发末期增加看门狗和测试用例；安全活动必须从 Item 定义、危害分析、安全目标、架构分配、硬件诊断、软件实现一直延伸到生产、维护和退役。

IEC~61508 给出跨行业的功能安全基础框架，道路车辆通常采用 ISO~26262，机械领域还会结合 ISO~12100、ISO~13849-1 或 IEC~62061，医疗设备则需要把 ISO~14971 风险管理与 IEC~62304 软件生命周期结合使用\cite{iec61508,iso26262,iso12100,iso13849-1,iec62061,iso14971,iec62304}。SIL、ASIL、PL 等等级来自不同风险模型，不能只按名称或数字直接换算。本章讨论可跨领域复用的工程方法，不替代目标行业的法规、标准解释和独立评估。

\section{安全、可靠性、信息安全与预期功能安全}

可靠性关注系统在规定条件和时间内持续完成要求功能的能力；功能安全关注由故障或失效行为引起的风险；信息安全关注恶意行为造成的机密性、完整性、可用性和控制权风险。三者相互影响，但证据不能互相替代。

一个系统可能经常检测到故障并安全停机，因此可用性不高但功能安全；也可能长期稳定运行，却在一个罕见传感器卡死时持续输出危险命令。网络攻击同样可能触发危害，因此安全分析需要与 \mbox{第~\ref{chap:security-lifecycle}~章} 的信息安全生命周期交换资产、攻击路径、更新和事件响应假设。

功能安全通常从“系统发生故障”出发；预期功能安全还关注系统没有传统故障、但感知能力不足、场景假设错误或功能边界不充分时产生的风险。ISO~21448 在道路车辆领域规定了 Safety of the Intended Functionality 的框架\cite{iso21448}。无论采用何种标准，团队都要明确当前论证覆盖故障、性能限制、误用、环境边界还是其中的组合。

\section{标准体系、Item 与安全生命周期}

安全生命周期开始于系统边界，而不是源码目录。Item 定义应列出功能、运行模式、用户、环境、能源、传感器、执行器、通信、维护接口、软件更新、硬件变型和外部依赖。若边界遗漏上位机、机械储能、反向供电或维修工具，后续危害分析会在错误系统上得到“完整”结果。

\begin{table}[H]
  \centering
  \caption{功能安全生命周期的主要阶段门}
  \begin{tabular}{p{25mm}p{39mm}p{43mm}p{29mm}}
    \hline
    阶段 & 核心问题 & 主要输出 & 退出证据 \\
    \hline
    概念 & 什么可能造成危害？ & Item、场景、危害与安全目标 & 风险分类获批 \\
    架构 & 谁承担何种安全职责？ & 需求分配、独立性与故障反应 & 架构分析闭合 \\
    实现 & 机制是否满足契约？ & 硬件、软件、工具与配置基线 & 单元和集成证据 \\
    确认 & 整机是否满足安全目标？ & 目标环境测试与安全确认 & 残余风险接受 \\
    运行 & 现场假设是否仍成立？ & 生产、维护、事件与变更记录 & 持续安全论证 \\
    \hline
  \end{tabular}
\end{table}

项目应建立适用性矩阵：每项法规或标准的版本、适用产品、责任人、裁剪理由和交付物是什么。仅声明“按 IEC~61508 开发”无法说明采用哪一安全完整性等级、哪些生命周期活动被裁剪以及谁接受残余风险。

硬件和软件生命周期可以采用不同过程，但必须在共同的 Item、接口、故障假设和配置基线上汇合。系统工程中的需求、接口和变更方法见 \mbox{第~\ref{chap:system-engineering-method}~章}。

\section{危害分析、风险分类与安全目标}

危害分析从系统功能、运行环境、人员暴露和可控性出发，识别危险事件并按适用标准分类。汽车 HARA 常组合严重度、暴露率和可控性，机械风险评估会考虑伤害严重度、发生概率和避免可能性；不同方法中的字母和等级不能直接比较。

分析方法应组合使用：

\begin{itemize}
  \item HAZOP、场景分析或操作分析用于发现偏差和危险事件；
  \item FMEA 从组件失效向上追踪局部影响和系统后果；
  \item FTA 从顶层危害反推必要故障组合和共同依赖；
  \item 时序与状态分析检查正确组件在错误顺序、模式或时间下的危险交互；
  \item 现场事件、相似产品和维修数据用于挑战纸面故障清单。
\end{itemize}

安全目标描述必须避免的系统级危害状态，例如“检测到位置反馈不可信后，在规定时间内撤销驱动输出”，而不是直接指定某个函数、看门狗型号或通信 CRC。实现机制应在后续需求分配中选择，避免过早锁定一个无法处理共因失效的方案。

分析需要包含正常使用、合理可预见误用、启动、关机、校准、维护、更新、降级和恢复。只分析稳态运行会遗漏输出上电毛刺、维修模式绕过、升级后配置丢失和故障复位后自动重启等高风险状态。

\section{安全状态、降级状态与 FTTI}

系统必须为每类危险事件定义安全状态。安全状态可能是断电、制动、保持受控输出、限速运行、切换冗余通道或把控制权交给本地机械保护，并不总是“全部关闭”。同一设备在有人区、无人区和维修模式下也可能需要不同反应。

降级状态允许在能力受限时继续提供部分服务，但必须规定剩余功能、允许持续时间、告警、退出条件和再次故障时的动作。把所有故障都立即停机可能产生次生危险；把所有故障都归为“可降级”则会掩盖风险。

故障容忍时间间隔（Fault Tolerant Time Interval, FTTI）限定从故障发生到危害形成的最大时间。最坏任务相位、诊断采样、去抖确认、调度阻塞、通信、输出驱动和机械能量变化都要进入预算：

\begin{equation}
  T_{sample}+T_{confirm}+T_{schedule}+T_{communicate}
  +T_{actuate}+T_{physical}<T_{FTTI}
\end{equation}

若诊断要求连续三个错误样本，最坏检测时间不是一个采样周期。若故障恰好发生在一次检查之后，还要加入完整任务相位。预算应保留明确裕量，并在最低电压、最高温度、最大总线负载和最坏机械负载下实测关键项。

进入安全状态后的重新使能同样属于安全需求。瞬态故障可以自动恢复还是必须锁存、需要何种自检、是否要求物理确认和谁有授权，都应写入状态机；无限复位或无条件恢复最后输出不是安全策略。

\section{安全需求、分配与双向追踪}

安全目标进一步分解为功能安全需求、技术安全需求以及硬件和软件安全需求。每条需求必须单义、可验证，并声明适用模式、故障条件、时间、接口和最终状态。

\begin{table}[H]
  \centering
  \caption{安全需求记录的最小字段}
  \begin{tabular}{p{28mm}p{48mm}p{56mm}}
    \hline
    字段 & 需要回答的问题 & 常见缺陷 \\
    \hline
    来源 & 对应哪个危害、安全目标或上级需求？ & 没有风险依据 \\
    条件 & 在什么模式、输入和故障下生效？ & 只写正常路径 \\
    义务 & 谁在多长时间内执行何种动作？ & “及时”“尽量”不可测 \\
    接口 & 数据、单位、范围、时序和错误是什么？ & 隐含跨模块假设 \\
    验证 & 用分析、测试、检查还是演示关闭？ & 测试无法触达 \\
    状态 & 草案、批准、实现、验证和偏差如何管理？ & 旧需求继续生效 \\
    \hline
  \end{tabular}
\end{table}

追踪必须双向成立：从危害能够找到安全目标、架构、代码和测试；从安全机制也能找到其需求来源。没有上级需求的“孤儿安全代码”可能是隐藏假设，没有验证证据的需求则仍然开放。CI 可以检查稳定需求 ID、测试结果和发布 manifest 的关联，但不能自动判断需求语义是否正确。

跨硬件与软件的需求要明确责任边界。例如“驱动应在 $20\,ms$ 内关闭”必须拆成软件请求时间、总线时延、门极关断、电流衰减和回读确认，而不能由多个团队各自假定其余部分时间为零。

\section{安全架构、独立性与免受干扰}

常见安全机制包括双通道比较、不同原理冗余、锁步核、ECC/Parity、独立看门狗、电压与时钟监控、端到端通信保护和安全岛。冗余只有在共同依赖和失效相关性受到控制时才有效；两个相同传感器共享同一电源、连接器和算法，不能自动视为独立通道。

安全监控器应比被监控功能简单，并具有直接进入安全状态的能力。复杂 Linux 应用可以提出控制请求，但最终输出约束可由 MCU、安全岛、独立逻辑或机械保护执行。监控器如果依赖同一调度器、堆、时钟和电源，就可能与被监控功能同时失效。

不同完整性等级的软件共享平台时，需要证明 freedom from interference：

\begin{itemize}
  \item 空间：内存、MMIO、DMA 和外设所有权不能越界；
  \item 时间：CPU、IRQ、锁、总线和看门狗服务具有上界；
  \item 通信：消息身份、长度、序号、新鲜度和错误传播受控；
  \item 共享资源：时钟、电源、缓存、内存带宽和恢复动作不会形成隐含耦合；
  \item 开发与配置：构建选项、校准、调试和更新不能绕过安全分区。
\end{itemize}

安全分解不能只画两个方框。每个通道的故障检测、独立资源、共同原因和组合失效概率都要有证据。若无法证明独立，就应按单一通道评估，或增加物理和设计多样性。

\section{随机硬件失效、FMEDA 与诊断覆盖}

功能安全通常区分随机硬件失效与系统性失效。前者来自器件在运行期的随机故障，需要失效率、失效模式和诊断分析；后者来自需求、设计、实现、工具或流程缺陷，不能仅靠增加失效率数字处理。

FMEDA 在 FMEA 基础上为元件和失效模式分配失效率，判断其安全、危险已检测、危险未检测、单点、残余或多点影响，并关联具体安全机制。失效率数据必须说明来源、环境、任务剖面、温度、应力和置信范围；不能把一个供应商 FIT 数字无条件复制到不同器件、封装和工况。

对规定危险失效集合，诊断覆盖可写成：

\begin{equation}
  DC=\frac{\lambda_{\mathrm{DD}}}
  {\lambda_{\mathrm{DD}}+\lambda_{\mathrm{DU}}}
\end{equation}

其中 $\lambda_{\mathrm{DD}}$ 为危险且可检测的失效率，$\lambda_{\mathrm{DU}}$ 为危险且未检测的失效率。这个比例只有在失效分类、检测时间和反应路径明确时才有意义。一个上电自检若无法在 FTTI 内发现运行期故障，不能对所有运行时间宣称同等覆盖。

IEC~61508 的 PFD/PFH、ISO~26262 的 SPFM、LFM 和 PMHF 服务于各自场景和判据\cite{iec61508,iso26262}。项目应使用适用标准的计算方法和目标，不能把不同指标拼成一个“安全百分比”。

ECC 能检测或纠正规定数量的位错误，但仍需分析地址译码、控制逻辑、未覆盖存储、scrub 周期和不可纠正错误反应。诊断机制本身也可能失效，必须纳入分析。

\section{共因失效、潜伏故障与独立性论证}

共因失效会让多个看似冗余的通道同时失效。典型来源包括共同电源、时钟、连接器、布线区域、温度、软件需求、编译器、标定数据、维护操作和同一批次器件。设计多样性可以降低部分相关性，但不会自动消除共同规格错误。

Dependent Failure Analysis 应沿能源、时序、空间、环境、数据和组织关系检查故障传播。分析结果需要落到具体控制，例如独立欠压监控、物理分隔、不同量程传感器、独立输出切断、只读安全标定和交叉团队评审，而不是只写“通道彼此独立”。

潜伏故障在单独发生时可能不产生可见危害，却会削弱下一次故障的保护能力。例如备用切断晶体管已经开路，只有主开关随后短路时才暴露。系统应通过周期诊断、启动自检、在线 test pulse 或维护 proof test 控制潜伏时间。

诊断和 proof test 间隔必须来自多点故障分析、故障到达率、运行时长和可维护性。测试过于频繁会增加磨损和误动作，过于稀疏则扩大潜伏窗口；测试本身的输出扰动和失败恢复也要进入安全分析。

\section{安全通信与端到端保护}

底层 CAN、Ethernet、SPI、UART 或共享内存的链路 CRC 只能说明某一传输层检测到部分位错误，不能证明消息来自正确发送者、属于正确变量、按正确顺序到达或仍然新鲜。安全通信应在端点之间建立独立于底层“黑通道”的保护层。

典型安全消息包含数据 ID、发送方和接收方、序号、时间或 age、payload 长度、状态以及端到端 CRC/MAC。接收方还要检查重复、跳号、超时、冻结值和模式组合：

\noindent\begin{minipage}{\textwidth}
\begin{lstlisting}[language=C,caption={安全通信封装的逻辑字段（示意）}]
struct safety_frame {
    uint32_t data_id;
    uint16_t source_id;
    uint16_t destination_id;
    uint32_t sequence;
    uint32_t age_ticks;
    uint16_t payload_length;
    uint16_t status;
    uint8_t  payload[MAX_SAFETY_PAYLOAD];
    uint32_t end_to_end_crc;
};
\end{lstlisting}
\end{minipage}

CRC 多项式、保护长度和残余错误率应与消息长度、更新率和目标完整性匹配。序号回绕、发送方重启和接收方切换必须有协议语义；把任何非连续序号都当成永久故障，会使正常重启造成不可恢复停机。

通信超时必须从 FTTI 预算推导。应用收到“最后一个有效值”不表示数据仍可使用，缓存和网关需要保留原始生成时间和 safety generation。安全接口的 schema、兼容与幂等契约见 \mbox{第~\ref{chap:interface-protocol-data-contracts}~章}。

\section{安全软件架构与状态机}

安全软件应限制动态行为和隐式共享状态。推荐使用静态或有上界资源、明确任务周期、受控 IPC、范围和单位检查、单向依赖以及集中安全输出仲裁。动态分配并非绝对禁止，但必须证明耗尽、碎片、失败和恢复不会破坏安全需求。

输入可信度不能只靠单次范围检查。还应检查来源、CRC、序号、时间戳、变化率、关联信号和当前模式：

\noindent\begin{minipage}{\textwidth}
\begin{lstlisting}[language=C,caption={带新鲜度和合理性检查的安全输入}]
bool accept_sample(const struct sample *s, uint32_t now)
{
    if (!crc_valid(s) || s->source_id != EXPECTED_SOURCE)
        return false;
    if (sequence_distance(last_seq, s->sequence) != 1U)
        return false;
    if ((now - s->timestamp) > MAX_SAMPLE_AGE)
        return false;
    if (!in_range(s->value) || !rate_is_plausible(last_value, s->value))
        return false;
    last_seq = s->sequence;
    last_value = s->value;
    return true;
}
\end{lstlisting}
\end{minipage}

安全状态机应显式区分初始化、自检、待命、活动、降级、安全锁存和维护。每个转换要有 guard、动作、超时和输出不变量。普通业务状态机只能请求使能，不能直接绕过安全仲裁写执行器。

错误处理本身也必须可分析。断言适合在开发阶段暴露违约，量产系统则应根据需求进入复位、降级或安全状态，并保存有界诊断。日志和崩溃转储不得阻塞安全输出，也不能因存储满而改变反应路径。

\section{时间行为、看门狗与资源预算}

实时安全需求需要最坏上界，而不是平均延迟。分析应包含 WCET、抢占、锁阻塞、IRQ/softirq、DMA、缓存和总线竞争、任务相位以及故障路径。Linux 实时调度和最坏延迟验证方法见 \mbox{第~\ref{chap:realtime-linux}~章}。

看门狗喂狗点必须证明关键工作完成，而不只是循环仍在运行。更可靠的 heartbeat 应包含任务 generation、输入和输出时间戳、状态、错误计数和关键依赖；监督器验证这些值按期推进后才服务硬件 watchdog。

windowed watchdog 可以检测过早和过晚服务，但窗口、独立时钟和复位传播都需验证。若 CPU、看门狗和最终安全输出共用同一失效电源域，复位请求可能永远无法到达执行器。

资源预算也属于安全契约。安全任务应有栈、队列、CPU、内存和总线带宽上限，并规定超限反应。诊断风暴、重复故障日志或无限重试不能消耗掉安全机制本身需要的资源。

对于多核或混合关键性平台，应限制 stop-machine、共享锁、缓存抖动、内存带宽和跨核复位的影响。空间隔离不自动提供时间隔离，平均负载较低也不能证明故障时存在调度裕量。

\section{系统性失效、编码与防御性设计}

系统性失效需要通过需求评审、架构约束、编码规范、静态分析、独立验证和变更控制降低。一个在所有设备上稳定复现的错误需求，不会因为硬件冗余而变成随机故障。

安全软件应特别控制：

\begin{itemize}
  \item 未定义行为、整数溢出、非法移位、越界和对象生命周期；
  \item 单位、量程、标度、字节序和校准 schema；
  \item 并发所有权、原子性、内存序和中断共享状态；
  \item 默认值、错误返回、部分成功和重启后的持久状态；
  \item 浮点 NaN/Inf、定点饱和、滤波器状态和控制输出限幅；
  \item 维护、调试、工厂测试与生产模式之间的权限转换。
\end{itemize}

“通过 MISRA 检查”不等于满足安全目标。编码规则只能降低部分缺陷概率，不能替代正确需求、调度分析、硬件故障分析和目标环境测试。对底层 C/C++、原子与 MMIO 边界的进一步讨论见 \mbox{第~\ref{chap:realtime-language-memory-model}~章}。

防御性检查需要对应明确反应。若所有异常都返回同一错误但上层继续使用旧输出，检查只是记录问题；若每个轻微异常都立即硬复位，又可能形成复位风暴。策略应按可恢复性、FTTI 和危害后果分级。

\section{工具置信度、复用软件与复杂组件}

编译器、链接器、代码生成器、静态分析器、测试工具、校准工具和烧录工具若可能引入或漏检安全错误，需要评估工具影响和错误检测能力。措施可以是限定使用、输出复核、回归对比、双工具链验证或采用经认可的 qualification evidence；不能只因为工具“行业常用”就视为可信。

复用 RTOS、Bootloader、Linux、协议栈或第三方库时，应记录版本、配置、使用需求、已知缺陷、测试历史、未使用功能和安全包装边界。质量管理级组件可以置于安全分区之外，但必须证明它不能干扰安全功能，并为其失效定义检测和降级。

Safety Element out of Context（SEooC）依赖一组假设需求，例如最大诊断周期、外部电压监控、允许温度、故障反应和集成测试。集成方必须逐项证明假设在目标产品中成立；只保存供应商 safety manual 而不执行约束，不能继承其安全声明。

复杂 MCU、SoC、PMIC 和传感器的安全机制常受初始化顺序、周期 test、寄存器锁、错误引脚和故障处理约束。驱动应把 safety manual 要求变成可审查配置和启动自检，并验证 warm reset、低功耗恢复和在线更新不会恢复为非安全默认值。

\section{验证、确认、覆盖与故障注入}

验证确认实现满足规定需求，确认判断整机是否在真实使用环境中满足安全目标。二者都需要与开发活动保持规定独立性；独立程度由适用标准和完整性等级决定，不能由同一作者仅自我签字关闭关键证据。

测试应分层覆盖需求、边界、鲁棒性、接口、时序、资源耗尽和恢复。结构覆盖用于发现未执行代码和测试缺口，但 100\% 语句覆盖不证明需求完整。MC/DC 等判据是否需要采用，应由目标行业和安全等级决定。

故障注入需要从分析得到 fault list，而不是随机拔线。注入点至少覆盖传感器卡死、偏移和慢漂移，通信丢失、重复和延迟，RAM/Flash 错误，任务超时、死锁和时钟异常，电源、温度与执行器故障。每次用例要记录：

\begin{itemize}
  \item 故障模型、位置、注入时刻、持续时间和可重复 seed；
  \item 预期检测机制、最大检测时间和目标安全状态；
  \item 实际内部事件、外部安全输出、机械或电气响应；
  \item 故障锁存、复位原因、重新使能和残余数据；
  \item DUT 失败与工装、仪器或脚本基础设施错误的区别。
\end{itemize}

故障注入不能破坏 oracle。若通过调试器暂停 CPU 注入错误，也同时改变 watchdog 和总线时序，就不能把结果直接解释为现场故障。主机、软件在环、虚拟平台、目标板和 HIL 的分层方法见 \mbox{第~\ref{chap:test-automation-hil}~章}。

长时间零失败只能形成统计证据，不能证明不存在系统性最坏路径。严格目标应结合分析、结构覆盖、故障注入、最坏时序测量和独立评审。

\section{生产、校准、维护与 Proof Test}

量产安全不能只依赖研发样机。生产流程应确认正确硬件变型、软件和标定版本，读取安全配置锁定位，测试最终输出切断路径，并把设备序列号、板卡修订、镜像摘要、校准和 EOL 结果关联起来。

校准参数应有单位、范围、硬件适用性、schema、CRC/签名和安全默认值。工站写入成功后必须由设备回读并执行独立合理性测试；缺失或越界标定不能静默回退到可能危险的通用值。

不同产品变型可能使用相同 PCB 和软件却具有不同执行器功率、机械限位或 FTTI。构建、设备树、标定和物料的组合必须受控，避免把通过验证的安全配置部署到错误硬件。

现场维护应规定 proof test 内容、间隔、人员资格、工具版本和失败处置。测试被跳过、设备长期离线或维修后未恢复安全锁定时，系统应产生可见状态，不能继续宣称原有诊断覆盖。

现场故障、误动作和 near miss 应反馈到危害分析、FMEDA 和测试计划。安全生命周期直到产品退役才结束；停产器件替换、维修换板和安全更新都需要变更影响分析。

\section{Safety Case、假设与变更管理}

Safety case 用结构化论证连接 claim、argument、evidence 和 assumption。顶层 claim 不是“系统安全”，而应限定产品版本、运行环境、寿命、模式和安全目标。证据包括分析、测试、评审、工具报告、生产控制和现场监测。

假设必须像需求一样受控，例如环境温度、外部电源质量、维护间隔、上位机超时、机械止挡和操作人员反应。一个外部假设若没有责任方、接口和验证方法，只是把风险移出文档。

发布基线应冻结安全需求、软硬件版本、标定、工具、未关闭偏差和证据摘要。任何硬件替代、编译器升级、RTOS 配置变化、诊断阈值或通信协议修改都可能使既有论证失效。

变更分析要从影响图出发：修改了哪些需求、接口、时间预算、失效模式、独立性假设、生产测试和现场恢复路径。仅重跑原有测试不能发现安全目标或架构假设已经变化。

临时偏差应具有风险评估、补偿措施、适用批次、到期条件和批准人。长期存在而无人复核的 waiver 会成为未建模的安全配置。

\section{信息安全与功能安全协同}

攻击者可以伪造传感器、延迟消息、耗尽 CPU、修改标定或阻止更新，从而触发与随机故障相同的危害。安全分析应把这些控制路径交给信息安全威胁模型，并把可利用结果反馈到功能安全危害和 FTTI 分析。

两类目标也可能冲突。安全系统希望快速开放维修接口以恢复设备，信息安全要求强认证；安全更新需要验证和防回滚，功能安全又要求更新期间保持确定输出和可恢复配置。解决方案应在共同状态机中定义权限、时间和降级，而不是让两个团队分别覆盖同一接口。

可信启动和证明可以确认加载状态，但不能证明控制算法满足安全目标；安全监控器可以限制输出，却不能保护密钥和更新来源。可信执行与设备身份见 \mbox{第~\ref{chap:trusted-execution-identity}~章}。

信息安全事件响应应包含安全影响分类：哪些设备可能失去独立保护，是否需要限权、现场检查或停用，以及补丁和密钥轮换是否改变安全证据。反之，现场安全事件也可能揭示可被攻击者利用的诊断绕过和维护后门。

\section{工程实例：受控执行器}

考虑一个由位置传感器、控制器、功率驱动和机械机构组成的线性执行器。普通功能接收目标位置并驱动电机；潜在危害是传感器、通信、软件或功率器件故障导致执行器越过安全边界并夹伤人员。

\subsection{Item、危害与安全目标}

Item 边界包括位置传感器、控制 MCU、驱动级、电机、电源、机械止挡、上位机通信和维修接口。示例安全目标定义为：检测到位置不可置信、控制任务失活、命令过期或越界运动后，在 $100\,ms$ 内撤销危险驱动，并阻止未经确认的自动重新使能。

该目标进一步分解为：

\begin{itemize}
  \item 控制器以不大于 $10\,ms$ 的周期验证位置范围、变化率和数据新鲜度；
  \item 独立硬件 watchdog 在 $40\,ms$ 内检测任务或时钟失活；
  \item 驱动使能具有默认关闭下拉，并可由独立安全输出撤销；
  \item 上位机命令包含数据 ID、序号、超时和目标范围；
  \item 位置回读与电流方向必须和命令一致，异常时进入安全锁存；
  \item 故障后只允许经过授权确认和完整自检重新使能。
\end{itemize}

\subsection{架构与 FMEA 片段}

普通运动控制器只能提交期望扭矩和方向，独立安全监督器检查位置、限位、电流、命令新鲜度和任务 heartbeat，并控制最终 enable。机械止挡和上游能量切断用于处理单一功率开关短路。

\begin{table}[H]
  \centering
  \caption{受控执行器的简化 FMEA}
  \begin{tabular}{p{25mm}p{30mm}p{38mm}p{33mm}}
    \hline
    失效模式 & 局部影响 & 检测机制 & 安全反应 \\
    \hline
    位置值卡死 & 控制持续输出 & 变化率、命令与位置不一致 & 撤销驱动并锁存 \\
    通信丢失 & 目标不再更新 & 序号、age 与命令超时 & 受控停止 \\
    控制任务卡死 & 输出保持旧值 & 独立 watchdog 与任务窗口 & 硬件关闭使能 \\
    功率开关短路 & 软件关闭无效 & 电流和位置回读 & 上游能量切断 \\
    安全输出开路 & 无法主动关断 & 启动 test pulse 与回读 & 禁止进入活动态 \\
    \hline
  \end{tabular}
\end{table}

“功率开关短路”说明安全分析必须跨越软件边界：若唯一开关失效为导通，继续发送软件关闭命令没有诊断价值，需要独立能量切断或机械限制。安全输出开路则说明保护机制本身也要诊断。

\subsection{FTTI 与资源预算}

一种示例时间分配为：最坏采样相位 $10\,ms$、传感器诊断和确认 $20\,ms$、安全任务调度 $10\,ms$、通信和输出传播 $5\,ms$、驱动撤销 $5\,ms$、机械能量衰减 $30\,ms$，总计 $80\,ms$，相对于 $100\,ms$ FTTI 保留 $20\,ms$ 裕量。

预算必须在最坏任务相位、总线拥塞和机械负载下测量。若需要连续三个错误样本，或机械衰减随温度和磨损变慢，相应上界必须重新分配。

安全监督器还需要独立栈、固定队列深度和输出寄存器访问权限。记录日志、上传遥测和普通 UI 不得持有阻止安全任务执行的锁。

\subsection{软件安全状态机}

\noindent\begin{minipage}{\textwidth}
\begin{lstlisting}[language=C,caption={执行器安全监督状态机（伪代码）}]
switch (safety_state) {
case SAFE_DISABLED:
    force_drive_off();
    if (authorized_reset() && all_self_tests_pass())
        safety_state = SAFE_READY;
    break;
case SAFE_READY:
    if (enable_request() && inputs_plausible())
        safety_state = SAFE_ACTIVE;
    break;
case SAFE_ACTIVE:
    if (!inputs_plausible() || command_expired() || !task_window_ok()) {
        force_drive_off();
        latch_fault();
        safety_state = SAFE_DISABLED;
    }
    break;
}
\end{lstlisting}
\end{minipage}

普通应用只能请求使能，不能直接写功率输出。安全监督器拥有最终权限，并使故障锁存、授权复位和自检成为显式可审查策略。复位后必须先保持 enable 关闭，再初始化输入和诊断，避免短暂默认值触发运动。

\subsection{验证矩阵}

\begin{table}[H]
  \centering
  \caption{受控执行器的故障验证矩阵}
  \begin{tabular}{p{27mm}p{39mm}p{39mm}p{29mm}}
    \hline
    故障 & 注入方法 & 验收门槛 & 证据 \\
    \hline
    传感器卡死 & 冻结位置并保持命令变化 & $T_{safe}<100\,ms$，故障锁存 & 输入、enable 与位置波形 \\
    命令重放 & 重复旧序号和有效 CRC & 拒绝旧命令，不恢复运动 & 帧、状态与错误码 \\
    任务死循环 & 在控制路径阻塞执行 & watchdog 关闭硬件使能 & heartbeat 与复位原因 \\
    主开关短路 & 故障夹具保持输出导通 & 独立能量切断有效 & 电流和切断时刻 \\
    安全输出开路 & 断开监督器关断路径 & 自检失败且禁止活动态 & test pulse 与回读 \\
    \hline
  \end{tabular}
\end{table}

所有测试都要测量故障发生到外部安全输出生效的时间，并检查故障锁存、复位原因和禁止自动重启。实例的重点不是固定阈值，而是从危害、安全目标、FTTI、架构分配、实现到故障注入形成可追踪证据链。

\section{数字化验收矩阵与评审清单}

项目应在安全计划中实例化 $T_{FTTI}$、允许最大检测时间 $T_{detect}$、安全输出时间 $T_{safe}$、目标诊断覆盖 $DC_{target}$、proof test 间隔 $I_{proof}$ 和故障注入规模 $N_{fault}$。这些符号不能留在量产发布门中，每个产品变型和运行模式都要给出具体数值与环境条件。

\begin{table}[H]
  \centering
  \caption{功能安全发布验收矩阵}
  \begin{tabular}{p{28mm}p{40mm}p{41mm}p{27mm}}
    \hline
    验收项 & 量化门槛 & 测试或分析条件 & 留存证据 \\
    \hline
    需求追踪 & 已批准安全需求双向覆盖，关键孤儿项为零 & 全部软硬件变型与偏差 & 追踪报告 \\
    FTTI & 最大 $T_{safe}<T_{FTTI}$ 且满足裕量 & 最坏相位、负载、温度和电压 & 波形与时序分解 \\
    诊断覆盖 & 对声明故障集 $DC\geq DC_{target}$ & FMEDA、诊断间隔和反应有效 & 失效模式清单 \\
    潜伏故障 & proof test 间隔 $\leq I_{proof}$ & 在线诊断、维护和跳过场景 & 测试记录与设备状态 \\
    故障注入 & 完成 $N_{fault}$ 个计划用例且反应符合需求 & 软件在环、目标板与 HIL 分层 & seed、日志与外部波形 \\
    重新使能 & 锁存故障未经授权和自检时零次恢复 & 复位、掉电、弱网与重复命令 & 状态轨迹 \\
    量产一致性 & 安全配置、标定和切断路径逐台通过 & 所有料号、工站和软件版本 & EOL 与追溯记录 \\
    \hline
  \end{tabular}
\end{table}

发布评审还应确认：

\begin{itemize}
  \item Item 是否覆盖用户、环境、能量、维护、更新和全部产品变型？
  \item 每个危害是否对应明确的安全目标、FTTI、安全状态和残余风险？
  \item SIL、ASIL 或 PL 是否仅在其适用标准和风险模型内使用？
  \item 安全需求是否能够双向追踪到架构、实现、测试和发布证据？
  \item 冗余通道的共因失效、潜伏故障和独立性是否经过分析？
  \item freedom from interference 是否覆盖空间、时间、通信和共享资源？
  \item FMEDA 的失效率、失效分类、诊断覆盖和测试间隔是否有依据？
  \item 安全通信是否检查身份、序号、新鲜度、长度、超时和模式？
  \item watchdog、安全输出和能量切断是否独立于被监控故障？
  \item 工具、第三方软件、SEooC 和复杂器件假设是否在目标产品中关闭？
  \item 故障注入是否测量外部安全反应，而不只检查内部日志？
  \item 生产、校准、维护、proof test、更新和退役是否进入 safety case？
  \item 信息安全威胁和补丁是否完成安全影响分析？
  \item 所有发布门槛是否已填入具体数值并关联原始证据？
\end{itemize}
